\begin{proposition}[A compact torus fibration with three singular fibres]\hypertarget{tag-02WR}{}\label{night:P4.2}\label{stmt:P4.2}
\StatementRecord{night:P4.2}{proved}{02WR}{\SeriesRef{night:P4.1}}

Let $E=\C/(\Z+i\Z)$ and $A=E^2$. Glue a principal $A$-bundle
over $\Pp^1$ by
$a_{\rm out}=a_{\rm in}+(\ell(t),0)$, where
$\ell(t)=\log(t)/(2\pi i)\in E$. Over disjoint discs about
$0,1,\infty$, replace its product models by the degeneration of
Proposition~\ref{night:P4.1}, using, respectively, the markings
\[
 C_i\times B_i\longrightarrow A,\qquad
 (z,w)\longmapsto(w,z),\quad(iz,iz+w),\quad(z,iz+w),
\]
and the punctured-disc product trivialization from the canonical section
of $\OO(R)$. These transitions define a smooth compact complex
threefold $f:Y\to\Pp^1$ with connected fibres, exactly three singular
fibres of the form in Proposition~\ref{night:P4.1}, and identity
monodromy on $H_1$ of the regular fibre.
\end{proposition}
\begin{proof}

\begin{enumerate}[label={$\langle1\rangle$\arabic*.},leftmargin=*]
\item The map $\ell(t)=(\log t)/(2\pi i)$ with values in $E$ is single-valued and
holomorphic on an annulus: changing the logarithm adds an integer.
Consequently $a_{\rm out}=a_{\rm in}+(\ell(t),0)$ is a holomorphic transition
for a principal $A$-bundle on $\Pp^1$. Product charts with compact fibre
prove that this bundle is Hausdorff and compact.

\item We choose disjoint discs about $0,1,\infty$ and replace their product models
by Proposition~\ref{night:P4.1} using the holomorphic group isomorphisms
\[
 C_i\times B_i\longrightarrow A:\qquad
 (z,w)\mapsto(w,z),\quad(iz,iz+w),\quad(z,iz+w),
\]
respectively. Multiplication by $i$ preserves $\Z+i\Z$. The three inverse maps
are $(x,y)\mapsto(y,x)$, $(x,y)\mapsto(-ix,y-x)$ and
$(x,y)\mapsto(x,y-ix)$, respectively. Their Gaussian-integer matrices
have determinants $-1,i,1$ and hence their real lattice determinants
are all one. These maps, $\lambda_B=v\lambda_A$, the canonical
punctured-disc trivialization in Proposition~\ref{night:P4.1}, and
the clutch $\ell$ specify all transitions.
Different base points separate in the base. Points over the same base
point lie in one Hausdorff local model, which separates them. The gluing is
therefore Hausdorff. Properness holds over each base chart and hence over
the base; compactness follows.

\item The two central components have connected elliptic fibres over connected
curves and meet each other. All remaining fibres equal $A$. The equations
$uv=t$ give exactly the indicated critical values. The punctured-disc
trivializations identify all fibre lattices with $\Lambda$, and the remaining
transition is a translation on $A$. Translation induces the identity on
$H_1(A;\Z)$, giving the asserted monodromy representation.

\end{enumerate}
\end{proof}